\section{Assessment}
\label{sec:assessment}

\subsection{What the formalization shows}

The repository proves, with no axioms beyond the classical base, that the two mechanisms in
the source papers actually produce what they claim: a forced Navier--Stokes flow that stays
in bounded energy while its velocity blows up, and an unforced Euler flow that leaves every
Sobolev class in finite time. The statements are the reference statements of the Millennium
problem, not paraphrases, and the two-dimensional structure of the argument---an explicit
self-similar core with an annular stress that is realised by oscillating pulses, and a
packet iteration that transfers a shear from one generation to the next---is visible in the
module structure rather than buried inside single monolithic proofs.

For a mathematician the interesting content is in the details that a paper leaves to the
reader and a proof assistant does not. Three examples from this development:

\begin{itemize}
\item \emph{Flatness is preserved, not re-derived.} The claim that the residual vanishes to
      every order at the singular point survives four different operations in the correction
      cycle. In a paper each operation gets a sentence (``the cutoffs are chosen so that the
      tails are smaller than any prescribed power''); in the formalization there is one
      predicate and a preservation lemma per operation, and the preservation lemmas are
      where the logarithmic weights come from.
\item \emph{The exponent ledger is arithmetic.} $\sigma_j=1/5+j/10$ is the whole content of
      the convergence of the correction scheme, and it is checked as a statement about real
      numbers in a file that knows nothing about fluid mechanics.
\item \emph{The pressure is never free.} In the Euler development the pressure is not stored
      as a field but as a gradient in a closed subspace of $L^2$, and the scalar pressure of
      the reference statement is reconstructed at the end. This is a stronger discipline than
      the papers' ``$p$ is determined up to a constant'', and it is what makes the
      reconstruction of the classical solution from the weak one checkable.
\end{itemize}

\subsection{What it does not show}

A formalization theorem is a statement about Lean terms, and the reader should keep in mind
what lies outside it.

\begin{itemize}
\item \emph{No new mathematics.} Every result here is a theorem of the source papers. The
      contribution is verification, not discovery.
\item \emph{Readability is not verified.} A proof term that typechecks can still be
      unreadable, and much of this development is machine-generated in style: long tactic
      blocks whose intermediate goals are only visible when replayed. The module docstrings
      and the statement layer are the parts written for humans.
\item \emph{Statement fidelity is an argument, not a theorem.} We argued in
      Section~\ref{sec:fidelity} that the statements say what the papers say. No formal
      system can check that argument; only a referee can.
\item \emph{Performance is part of correctness in practice.} A file that takes an hour to
      elaborate is a file nobody will modify. The largest file in this repository,
      \file{Euler/EulerProof.lean}, is 20{,}755 lines in one module, and it is the place
      where a future change is most expensive.
\end{itemize}

\subsection{What the migration shows}

The migration from \Lean{} \texttt{v4.34.0-rc2} to \texttt{v4.36.0-rc1} is a controlled
experiment on how much of a large development depends on the exact shape of a library API.
Our observations:

\emph{The breakage is concentrated in the analysis API, not in the proofs.} The failures we
found are overwhelmingly of three kinds---a structure became a definition, a hypothesis moved
from implicit to explicit, an argument was reordered---and each is repaired locally, in a
handful of lines. What made them expensive is not their difficulty but their number and the
fact that each requires a full re-elaboration of a file to test.

\emph{Statement changes propagate.} When \lean{eLpNorm} changed its value on
non-measurable functions from $0$ to $\infty$, a project lemma
(\lean{lpNorm\_mono\_of\_norm\_le\_ae}) became false without an extra measurability
hypothesis. A proof repair would have hidden this; the correct response was to strengthen the
hypothesis and let the call sites fail. Formalization projects should expect a small number of
such semantic changes per library generation, and should budget for them.

\emph{Downstream libraries pin harder than \Mathlib{}.} \TauCeti{} follows \Mathlib{}
\texttt{master} and pins its own toolchain; Comparator is released per toolchain. A project
that wants all three is forced onto the intersection of their toolchains, which in our case
excluded the newest \Mathlib{} release candidate and excluded Comparator altogether. The
ecosystem would benefit from a released ``analysis layer'' that tracks a \Mathlib{} release
cycle rather than its development branch.

\subsection{Recommendations}

For projects of this kind we would recommend four things, each of which we would have wanted
here.

\begin{enumerate}
\item \emph{Write the statement layer first and keep it small.} A reader, and a referee,
      should be able to check the meaning of the theorem without reading the proof. The
      module \file{NavierStokes/R3/ProblemStatement.lean} is a good model: it contains the
      theorem as a named proposition, states that nothing in the module proves it, and lets
      the proof live elsewhere under the same name.
\item \emph{Quarantine arithmetic bookkeeping.} Exponent ledgers, smallness hierarchies and
      ``choose $\varepsilon$ so that $C\varepsilon^a<1$'' chains should be their own modules
      about real numbers. They are where errors in a paper hide, and they are cheap to check.
\item \emph{Adopt the module system and pin full commit hashes.} The module system makes
      import hygiene checkable and removes a whole class of warnings; full hashes make the
      build reproducible and the migration a decision rather than an accident.
\item \emph{Record the trust base in the repository.} \lean{\#print axioms} next to the
      headline theorems, a list of intentional \lean{sorry}s, and a note on which parts are
      conditional. The cost is a few lines; the benefit is that the claim ``this is verified''
      becomes checkable by a reader who does not trust the authors.
\end{enumerate}
